% SHARD-ID: CA-45-QUBIT-SYMMETRY-EXCLUSION-ROADMAP
% SHARD-TITLE: Roadmap for Qubit Symmetry Exclusion
% SHARD-SUMMARY: Summarizes the necessary-equation block and turns it into an implementation plan for algebraic filters and SDP exclusions.
% SHARD-SUMMARY: Records the delegated subagent work, the checked Julia surface, and the open convention gaps.
% SHARD-SUMMARY: Gives acceptance tests for future runs before any Hamiltonian is claimed excluded.
% SHARD-KEYWORDS: roadmap, qubit exclusion, SDP, Poincare, Virasoro, orchestration

\section{Roadmap for Qubit Symmetry Exclusion}
\label{sec:qubit-symmetry-exclusion-roadmap}

\subsection{Status, Sources, and Dependencies}

\textbf{Status: implementation plan.}  This closes the current qubit
necessary-equation expansion and names the next build steps.

Dependencies: CA-34--CA-44, CONVENTIONS.md (m)--(o).

% Review orchestration: reviews/2026-05-31_qubit_nn_sdp_hierarchy/ORCHESTRATION.md.

\subsection{What We Have Now}

The one-dimensional first-moment route now has explicit necessary equations on
the Pauli coefficient matrix \(h_{\alpha\beta}\):
\[
  p_{aed}
  =
  -2\sum_{b,c=1}^3h_{ab}h_{cd}\epsilon_{bce},
\]
\[
  A(h)=D u,
  \qquad
  B(h)-u-v^2(h-eI)=D w.
\]
These equations already rule out fully local symmetric on-site Hamiltonians and
commuting classical densities for any nontrivial Poincare route of this form,
because their generated momentum current is zero.

The two-dimensional shard gives a square-lattice edge schema, not yet code.
The Witt/Virasoro shard gives exact finite mode commutators and residuals, but
keeps Koo-Saleur's finite-nonclosure warning active.  The vacuum/SDP shards turn
``there exists a global state \(\omega\)'' into finite positive moment data and
linear relation constraints.

\subsection{Implementation Steps}

\begin{enumerate}
\item Add symbolic coefficient builders for \(A(h)\), \(B(h)\), and
  coboundary equations, returning sparse polynomial tensors rather than dense
  matrices.
\item Add a solver-facing algebraic filter: for a fixed \(h\), check current
  collapse; optionally solve the linear coboundary witnesses \(u,w\) after
  \(p,A,B\) are computed.
\item Build a Pauli-word moment enumerator for intervals \([-N,N]\), including
  product constants, translation-invariance canonicalisation, and Hermitian
  realification.
\item Export the fixed-\(h\) feasibility problem (44.1)--(44.4) through a local
  SDP interface.  Record every run under \texttt{runs/} with the command line,
  \(h\), level \(N\), parameters \(e,v^2,u,w\), and solver status.
\item Add mutation tests: perturb a known zero-current density to make the
  current nonzero; perturb the coboundary sign to make the coboundary
  reconstruction test fail; perturb a residual coefficient and require a small
  SDP level to change.
\end{enumerate}

\subsection{Acceptance for an Exclusion Claim}

A future statement ``\(h\) is excluded at level \(N\)'' must point to:

\begin{itemize}
\item the exact Pauli coefficient matrix \(h_{\alpha\beta}\) and density
  convention;
\item the algebraic residual tensors and any witnesses \(e,v^2,u,w\);
\item the finite word sets \(\mathcal W_N,\mathcal A_N\);
\item the SDP solver transcript or certificate in a run bundle;
\item a worklog entry explaining why the infeasible constraints are necessary
  for the named route.
\end{itemize}

\subsection{Open Gaps}

The central charge convention remains unfixed.  The two-dimensional edge schema
has no implementation.  The SDP hierarchy is fixed-\(h\); scanning continuous
families of \(h_{\alpha\beta}\) requires a separate polynomial relaxation.
None of the filters replaces CA-15's scaling-limit obligations.  They are
designed to cheaply reject impossible candidates before the expensive continuum
work begins.

